% Options for packages loaded elsewhere
\PassOptionsToPackage{unicode}{hyperref}
\PassOptionsToPackage{hyphens}{url}
\PassOptionsToPackage{dvipsnames,svgnames,x11names}{xcolor}
\documentclass[
  11pt,
]{article}
\usepackage{xcolor}
\usepackage[margin=2.4cm]{geometry}
\usepackage{amsmath,amssymb}
\setcounter{secnumdepth}{-\maxdimen} % remove section numbering
\usepackage{iftex}
\ifPDFTeX
  \usepackage[T1]{fontenc}
  \usepackage[utf8]{inputenc}
  \usepackage{textcomp} % provide euro and other symbols
\else % if luatex or xetex
  \usepackage{unicode-math} % this also loads fontspec
  \defaultfontfeatures{Scale=MatchLowercase}
  \defaultfontfeatures[\rmfamily]{Ligatures=TeX,Scale=1}
\fi
\usepackage{lmodern}
\ifPDFTeX\else
  % xetex/luatex font selection
\fi
% Use upquote if available, for straight quotes in verbatim environments
\IfFileExists{upquote.sty}{\usepackage{upquote}}{}
\IfFileExists{microtype.sty}{% use microtype if available
  \usepackage[]{microtype}
  \UseMicrotypeSet[protrusion]{basicmath} % disable protrusion for tt fonts
}{}
\usepackage{setspace}
\makeatletter
\@ifundefined{KOMAClassName}{% if non-KOMA class
  \IfFileExists{parskip.sty}{%
    \usepackage{parskip}
  }{% else
    \setlength{\parindent}{0pt}
    \setlength{\parskip}{6pt plus 2pt minus 1pt}}
}{% if KOMA class
  \KOMAoptions{parskip=half}}
\makeatother
\setlength{\emergencystretch}{3em} % prevent overfull lines
\providecommand{\tightlist}{%
  \setlength{\itemsep}{0pt}\setlength{\parskip}{0pt}}
% v3 header: single column, tight reference list, clickable links.
% Loaded via pandoc -H. Fixes for the six issues Havid found in the v2 PDF are
% mostly in the markdown generator; this file only carries LaTeX-side settings.
\usepackage{graphicx}
\usepackage{float}
\usepackage{caption}

% Figure captions: LaTeX owns the "Figure N." label. The v2 PDF rendered
% "Figure 4: Figure 4." because the caption text ALSO began with "Figure 4.",
% written into the markdown by the generator. v3 passes a bare caption and lets
% LaTeX number it, so the label appears exactly once.
\captionsetup[figure]{labelfont=bf, labelsep=period, font=small,
                      justification=justified, singlelinecheck=false,
                      skip=6pt}

% Tighten the numbered reference list produced by pandoc's default
% CSLReferences environment.
\usepackage{etoolbox}
\AtBeginEnvironment{CSLReferences}{%
  \small
  \setlength{\parskip}{0pt}%
  \setlength{\itemsep}{0pt}%
  \setlength{\topsep}{2pt}%
}

% No section numbering from LaTeX: the generator already writes "1. Introduction",
% "2. The Certified Frontier ...", so LaTeX numbering would double it.
\setcounter{secnumdepth}{0}
\usepackage{bookmark}
\IfFileExists{xurl.sty}{\usepackage{xurl}}{} % add URL line breaks if available
\urlstyle{same}
% fallback for those not using the hyperref driver hyperxmp:
\makeatletter
\@ifundefined{xmpquote}{\newcommand{\xmpquote}[1]{#1}}{}
\makeatother
\hypersetup{
  colorlinks=true,
  linkcolor={[HTML]{1A4E8A}},
  filecolor={Maroon},
  citecolor={Blue},
  urlcolor={[HTML]{1A4E8A}},
  pdfcreator={LaTeX via pandoc}}

\author{}
\date{}

\begin{document}

\setstretch{1.05}
\begin{center}
{\LARGE\bfseries Perovskite Photovoltaics in July 2026: Buried-Interface Chemistry, Defect Tolerance and the Certified Efficiency Frontier\par}
\vspace{1.1em}
{\large Havid Aqoma\textsuperscript{1,2,3,4}\par}
\vspace{0.7em}
\begin{minipage}{\textwidth}\centering\small
\textsuperscript{1}School of Energy and Chemical Engineering, Xiamen University Malaysia, Selangor Darul Ehsan 43900, Malaysia\\[0.25em]
\textsuperscript{2}Kelip-kelip! Center of Excellence for Light Enabling Technologies, School of Energy and Chemical Engineering, Xiamen University Malaysia, Selangor Darul Ehsan 43900, Malaysia\\[0.25em]
\textsuperscript{3}Centre of Excellence for Green and Advanced Technologies in Efficient Separations (GATES), School of Energy and Chemical Engineering, Xiamen University Malaysia, Selangor Darul Ehsan 43900, Malaysia\\[0.25em]
\textsuperscript{4}Gulei Innovation Institute, Xiamen University, Zhangzhou 363200, China\\[0.25em]
\end{minipage}
\vspace{2.2em}
{\small September 8, 2026\par}
\end{center}
\vspace{1.0em}

\subsection{Abstract}\label{abstract}

Perovskite photovoltaics published 648 indexed works in July 2026, of
which 181 were read closely and 80 are cited here. This review is
organised mechanistically rather than chronologically, and separates
what the month's literature verified from what it claimed. The certified
frontier advanced on three fronts, each through interfacial rather than
bulk chemistry. Noncoplanar multidentate contacts reinforcing the buried
interface delivered a 33.1\% certified perovskite/silicon tandem
{[}1{]}; hydrogen-bonding-regulated crystallisation of the wide-bandgap
subcell reached 32.13\% certified {[}2{]}; and homogenisation of the
tin-lead distribution through A-site methylammonium cations gave 29.57\%
certified in an all-perovskite tandem {[}3{]}. Single-junction
certification clustered near 27.31\%, with a soft supermolecule
stabilising the buried interface at 27.12\% certified and carrying the
same chemistry to a 655.2 cm\(^2\) aperture module {[}4{]}, the largest
certified area in the month. Only 22 of 169 closely read papers reported
an independently certified power conversion efficiency at all.
Operational-stability evidence was the weakest link: 4 papers reported a
time to 80\% of initial performance (T80), the longest at 9800 h
{[}5{]}, and 6 named a protocol from the International Summit on Organic
Photovoltaic Stability (ISOS), so most lifetime claims this month cannot
be compared against each other. Across the corpus an efficiency value
was detected in at least 50.4\% of abstracts, independent certification
in at least 8.3\%, and a complete open-circuit voltage, short-circuit
current density and fill factor triplet in at least 2.9\%. Three
mechanistic conclusions follow. Buried-interface chemistry, not absorber
composition, now sets the achievable open-circuit voltage. Defect
tolerance has overtaken defect elimination as the practical design
strategy. Halide homogenisation in wide-bandgap absorbers is the
rate-limiting step for tandem integration, and it is the subcell least
often optimised in the month's work.

\textbf{Keywords:} perovskite photovoltaics; buried interface; defect
passivation; wide-bandgap perovskite; tandem solar cells; operational
stability

\subsection{1. Introduction}\label{introduction}

Perovskite single junctions are now certified above 27\% power
conversion efficiency (PCE) and perovskite/silicon tandems above 33\%,
so the field's open problems have moved from raw efficiency to whether a
claimed efficiency is verifiable, whether it survives at module area,
and whether it survives operation. High peak efficiencies on small-area
cells increasingly diverge from the operational demands of modules,
where non-radiative recombination losses at defect sites, interfacial
delamination, and volatile species desorption dictate practical
viability.

In July 2026, 648 indexed works met the scope gate and 181 were read
closely. This review follows four mechanistic threads: buried-interface
and contact chemistry, defect tolerance and ion migration, wide-bandgap
absorbers for tandems, and operational stability, with scale-up as the
constraint coupling them. Interfacial chemical interactions govern
carrier extraction and charge accumulation, defect landscapes dictate
electric-field-driven ion migration and halide segregation, and
wide-bandgap absorbers remain vulnerable to photoinduced phase
separation, with all these processes subject to the physical boundaries
of scalable deposition.

The month's headline numbers are strong, the evidence behind them is
unevenly reported, and this review separates the two. Critical
evaluation distinguishes verifiable physical mechanisms from uncertified
champion values, non-standardized aging procedures, and uncharacterized
degradation modes.

\subsection{2. The Certified Frontier in July
2026}\label{the-certified-frontier-in-july-2026}

Only 22 of 169 closely read papers reported an independently certified
power conversion efficiency (PCE), revealing that most literature relies
on unverified laboratory metrics. Certified performance ceilings diverge
across device architectures. For perovskite/silicon tandems, verified
efficiency reached 33.1\% in an inverted p-type-intrinsic-n-type (p-i-n)
structure {[}1{]} alongside 32.13\% in two-terminal (2T) devices
{[}2{]}. In all-perovskite 2T tandems, narrow-bandgap subcell
limitations and non-radiative recombination constrained the certified
ceiling to 29.57\% {[}3{]}, with lower certified tiers at 28.84\%
{[}6{]} and 27.49\% {[}7{]}. Single-junction p-i-n perovskite solar
cells (PSCs) converged between 27.3\% and 27.31\% {[}8,9,10{]}.

Device area exposes a fundamental evidentiary divide: a self-reported
champion on a sub-0.1 cm2 cell and a certified value on a module
represent distinct classes of physical evidence. Uncertified champion
micro-cells suppress lateral series resistance and avoid macroscopic
defects, allowing self-reported efficiencies to run ahead of certified
limits. Conversely, large-area modules must overcome interconnect
resistance, transparent conductor sheet resistance, and non-uniform film
crystallization. Demonstrating scalable transport, certified p-i-n
modules verified 27.31\% across both 766 cm2 {[}9{]} and 684.75 cm2
{[}8{]}.

Beyond these verified benchmarks, uncertified claims overstate
large-area performance. Uncertified reports claimed 27.12\% on 655.2 cm2
{[}4{]} and 26.53\% on 818 cm2 {[}11{]}, while self-reported in-lab
testing on a 625 cm2 module claimed 27.8\% despite achieving an
independent certification of only 27.3\% {[}10{]}. This discrepancy
highlights how uncertified testing masks transient capacitive currents
and aperture misalignment, confirming that unverified claims on sub-0.1
cm2 cells distort the true technological frontier.

\begin{figure}[htbp]
\centering
\includegraphics[width=0.95\textwidth]{fig/F1_certified_frontier.pdf}
\caption{The certified frontier in July 2026. Filled markers are independently certified efficiencies, open markers are self-reported champion values, grouped by device family. Horizontal bars mark the highest certified value in each family.}
\end{figure}

\subsection{3. Buried Interfaces and Contact
Chemistry}\label{buried-interfaces-and-contact-chemistry}

The buried interface between the hole-transport layer (HTL) and
perovskite absorber dictates carrier extraction kinetics and mechanical
reliability in inverted perovskite solar cells (PSCs) {[}12{]}. Weak
chemisorption and disordered self-assembled monolayer (SAM) packing on
metal oxides encourage molecular desorption and interfacial
delamination; removing loosely bound molecules increases the interfacial
fracture energy from 0.94 to 4.82 J m-2 {[}13{]}. To prevent interfacial
detachment, multidentate chelating architectures substitute weak
monodentate anchoring {[}8,14,15{]}. Tetraphosphonic acid iron(III)
porphyrins establish tetradentate chelation to oxide substrates,
yielding a certified power conversion efficiency (PCE) of 26.82\%
(27.06\% champion) {[}14{]}. Bisphosphonate head groups similarly
reconstruct nickel oxide (NiOx) interfaces by coordinating
undercoordinated lead (Pb2+) to achieve a certified 27.31\% PCE {[}8{]},
or reinforce anchoring via hydrogen bonding to retain over 90\% of
initial performance after 200 h of maximum power point tracking (MPPT)
{[}15{]}. Tetraphosphorylated phthalocyanines lock interfacial packing
to secure a certified 23.1\% PCE (22.9\% stabilised) on a 20.3 cm2
module with zero degradation across 61 days of outdoor testing under
International Summit on Organic Photovoltaic Stability (ISOS) ISOS-O-1
protocols {[}16{]}. Beyond anchoring strength, self-aggregation via
pi-pi stacking creates pinholes that expose bare transparent conductive
oxides. Introducing noncoplanar geometry suppresses aggregation to
achieve a certified 33.1\% tandem PCE (33.3\% champion; 23.7\%
single-junction) {[}1{]}. Alternatively, co-adsorbing Lindqvist
polyoxometalates increases the SAM anchoring ratio from 29.93\% to
53.98\% while reducing the desorption rate from 30.7\% to 5.6\% during
continuous MPPT {[}17{]}. Interfacial aggregation is also counteracted
using ionic liquid lubricants that preserve 92.3\% initial efficiency
after 800 h at 85 °C under air mass 1.5 global (AM 1.5G) illumination
{[}18{]}, proton-transfer complexes resisting 75 °C degradation for 1368
h {[}19{]}, or thermally activated in situ polymerization networks that
achieve an uncertified 26.87\% PCE {[}20{]}.

At the absorber contact, asymmetric nucleation generates buried
microvoids, tensile lattice strain, and deep defect states that induce
non-radiative recombination {[}21,22{]}. Bidentate phenanthroline
molecules chelate Pb2+ and facilitate dimethyl sulfoxide (DMSO) solvent
evaporation to eliminate interfacial voids, delivering a 25.15\% PCE,
but this evidence is constrained to an uncertified 0.04 cm2 aperture
area {[}22{]}. Supramolecular calixarenes mitigate defect propagation
via host-guest interactions, yielding a certified 27.12\% cell PCE
alongside a certified 22.25\% on a 655.2 cm2 module {[}4{]}.
Electrostatic alignment and defect passivation proceed concurrently
through molecular dipoles: sulfonate-anchored dipotassium salts
coordinate undercoordinated tin (Sn4+) sites on tin oxide (SnO2)
electron-transport layers (ETLs) to passivate oxygen vacancies {[}23{]},
while fluorinated carbazole self-assembled monolayers (SAMs) passivate
halide vacancies, scaling to 1 cm2 with 25.06\% PCE and 3.11\% relative
loss {[}24{]}. Polycyclic aromatic SAMs with tailored electrostatic
potentials (26.20\% PCE) {[}25{]}, multifunctional modifiers (26.67\%
PCE) {[}26{]}, ternary phosphine bridges stabilizing formamidinium
cations {[}27{]}, and preburied two-dimensional (2D) perovskite inducers
that invert tensile to compressive strain {[}21{]} similarly regulate
the heterojunction. Nevertheless, theoretical one-dimensional Solar Cell
Capacitance Simulator (SCAPS-1D) simulations predict ideal efficiencies
exceeding 30\%, including 30.81\% at 300 K {[}28{]} and 29.1\% in
HTL-free bilayer ETLs {[}29{]}, yet these drift-diffusion calculations
omit the spatial heterogeneity of SAM desorption, localized pinholes,
and void formation that dominate physical devices.

\subsection{4. Defect Tolerance, Passivation and Ion
Migration}\label{defect-tolerance-passivation-and-ion-migration}

Non-radiative recombination at carrier-selective contacts in perovskite
solar cells (PSCs) is driven by undercoordinated lead(II) (Pb2+)
cations, dangling bonds, and halide vacancies that require Lewis
acid-base coordination {[}30,31,32,33,34{]}. At hole-transport layer
(HTL) contacts, carbazole-based self-assembled monolayers (SAMs)
coordinate Pb2+ to reach 26.85\% power conversion efficiency (PCE)
{[}30{]}, while carboxyl-sulfur molecules coordinate dangling bonds and
free lead(II) iodide (PbI2) for an uncertified 26.7\% PCE {[}32{]}. At
electron-transport layer (ETL) contacts,
1-(4'-sulfophenyl)-3-carboxyl-5-pyrazolone sodium salt (SCP) yields a
certified 26.84\% PCE (champion 27.25\%, stabilised 26.84\%) with 90\%
retention after 1485 h of maximum power point tracking (MPPT) {[}31{]}.
Amphoteric ionic liquids passivate both acidic and basic sites for
26.13\% champion PCE (20.06\% PCE and 85.74\% fill factor (FF) in 1.84
eV cells) {[}33{]}, while Pb-N coordination and vacancy filling
preserves 88.67\% efficiency after 1000 h MPPT at 65 °C (uncertified
25.06\% PCE) {[}34{]}. Passivating tin(IV) oxide (SnO2) oxygen vacancies
with conjugated sorbate anions yields an uncertified 22.16\% PCE
(open-circuit voltage (Voc) 1.183 V, FF 84.30\%) {[}35{]}. Plasma
surface engineering passivates laser scribes, securing a certified
23.5\% PCE on a 65.05 cm2 module and 27.2\% small-cell champion PCE
(98.1\% retained after 2000 h MPPT at 85 °C) {[}36{]}.

Electric-field- and heat-induced migration of halides and A-site cations
accelerates phase segregation, requiring chemical trapping and
dimensional barriers {[}37,38,39,40{]}. Dual iodide-trapping molecules
immobilize mobile halides at the HTL interface for 26.81\% champion PCE
and a T93 lifetime over 1140 h at 85 °C {[}37{]}, while
acetylthiocholine chloride (ATC) coordinates Pb2+ and fills vacancies to
suppress segregation, yielding 22.68\% champion PCE on a 63.7 cm2 module
{[}40{]}. Two-dimensional/three-dimensional (2D/3D) heterojunctions
formed by halide imidazolium cations block ion pathways, retaining 92\%
efficiency after 1200 h at 85 °C (26.33\% champion PCE) {[}39{]}. In
situ two-dimensional (2D) capping via hydrogen bonds with formamidinium
(FA+) retains 90\% efficiency after 800 h at 85 °C in N2, yet evaluation
on a 0.072 cm2 aperture (26.32\% champion PCE) leaves scalability
unproven {[}41{]}. Molten salt dissolution of 2D phases suppresses
random nucleation and eliminates interfacial PbI2, boosting PCE from
23.43\% to 25.73\% {[}42{]}. Unified ligand-dimensional designs halt
cation migration, achieving certified 26.52\% PCE (26.88\% champion) on
small cells and certified 17.75\% PCE (20.25\% uncertified) on 900 cm2
modules, retaining over 91\% efficiency after 1000 h at 85 °C and 85\%
relative humidity (RH) {[}38{]}.

Residual tensile strain accelerates defect nucleation and ion transport,
coupling mechanical stress to electronic loss {[}43,44,45,46,47{]}.
Coordinating adjacent Pb2+ with dithiopyr relieves lattice strain and
directs preferential (100) orientation, yielding certified 26.14\% PCE
(26.9\% champion) and 21.28\% on 30 × 30 cm2 modules {[}43{]}.
Cross-linked 2-amino-3-hydroxybutanoic acid (AHBA) networks relax buried
SnO2 interface strain and accelerate electron extraction to 3.5 ps for
26.32\% champion PCE {[}44{]}. Low-temperature liquid-assisted
vibrational stress-relief (LTL-VSR) relieves macroscopic stress without
additives, yielding 26.47\% rigid and 24.64\% flexible cell
efficiencies, alongside 23.59\% (16.8 cm2) and 20.55\% (600 cm2) module
efficiencies {[}46{]}. Photoresponsive fullerene C60-azo generates
compressive stress under ultraviolet (UV) light via trans-to-cis
isomerization, retaining 92.7\% performance after 488 h UV exposure with
26.6\% champion PCE {[}45{]}. Finally, phenethylammonium iodide (PEAI)
paired with cerium oxide (CeOx) provides redox-buffering passivation,
retaining 95\% efficiency after 100 thermal shock cycles between 80 °C
and -80 °C (at 16 °C/min) with 25.3\% champion PCE in poly(triarylamine)
(PTAA)-based devices {[}47{]}.

\subsection{5. Wide-Bandgap Absorbers and Tandem
Integration}\label{wide-bandgap-absorbers-and-tandem-integration}

Monolithic tandems can surpass the \textasciitilde29.4\% crystalline
silicon (c-Si) efficiency limit {[}48{]}, yet two-terminal (2T)
integration remains constrained by halide phase segregation in
wide-bandgap (WBG) top absorbers and interfacial non-radiative
recombination. In-operando thermal cycling between -100 and +100 °C
under air mass 0 (AM0) illumination reveals that WBG subcells segregate
rapidly above 50 °C with only partial recovery during dark cooling,
accelerating non-radiative voltage loss {[}49{]}. Suppressing this
demixing requires reinforcing octahedral framework integrity and
passivating grain boundaries {[}50,51,52{]}. B-site strontium (Sr2+)
alloying strengthens metal-halide bonding and weakens electron-phonon
coupling, yielding uncertified 31.66\% power conversion efficiency (PCE)
in fully-textured perovskite/silicon tandems with an operational T80
lifetime (time to 80\% initial efficiency) of 6,785 hours {[}50{]}. At
grain boundaries, TPB-Br molecular extrusion drives halogen
homogenization to deliver 23.63\% PCE in 1.68 electronvolt (eV) subcells
and an uncertified 32.03\% on 1.05 cm2 tandems {[}51{]}, while
conductive ionic-liquid/polyurethane polymer bridges passivate defect
states to yield uncertified 20.85\% WBG and 25.92\% flexible tandem PCEs
{[}52{]}. Ionic-liquid additives similarly confer radiation resilience,
yielding certified 27.49\% monolithic tandems over 12.56 cm2 that
maintain 93\% efficiency under proton irradiation {[}7{]}.

Defective buried contacts and unguided nucleation introduce large
open-circuit voltage (Voc) deficits. Retarding nucleation via
coordinating solvents directs micrometer-scale grain growth {[}53,54{]}:
gamma-valerolactone (GVL) coordination with formamidinium (FA+) cations
achieves uncertified 32.5\% champion tandem PCEs {[}53{]}, while
GVL-dimethyl sulfoxide (DMSO) inks yield printed cells exceeding 17\%
and 28\% tandems {[}54{]}. Regulating crystallization and passivating
defects at hole-transport layers (HTLs) further limits interface
recombination {[}2,55,56{]}. Inserting a 3-TMA interlayer on
self-assembled monolayers (SAMs) regulates crystallization through
hydrogen bonding, achieving 23.17\% (1.67 eV) and 19.61\% (1.84 eV)
subcells and a certified 32.13\% (33.21\% champion) silicon tandem
{[}2{]}. Bilayer graphene oxide/self-assembled monolayer (SAM) HTLs
similarly stabilize triple-junction devices at 27.0\% stabilized PCE
(27.3\% champion) {[}55{]}, whereas omitting nickel oxide (NiOx) HTLs
via hybrid deposition yields flexible perovskite/copper indium gallium
selenide (CIGS) tandems with certified 25.9\% PCE (26.4\% champion, 1.85
V Voc) {[}56{]}.

\begin{figure}[htbp]
\centering
\includegraphics[width=0.95\textwidth]{fig/F3_effort_map.pdf}
\caption{Where the month's mechanistic effort was directed. Depth-tier papers by mechanism axis, stacked by device family.}
\end{figure}

\subsection{6. Operational Stability and Degradation
Pathways}\label{operational-stability-and-degradation-pathways}

Depth-resolved profiling reveals that operational photo-degradation
initiates at both buried and top interfaces before propagating into bulk
perovskite absorbers {[}57{]}. At buried hole-transport layer (HTL)
contacts in inverted perovskite solar cells (PSCs), molecular disorder
accelerates non-radiative recombination. Rigidifying molecular
orientation by coassembling carbazole self-assembled monolayers (SAMs)
suppresses interfacial defects, yielding a certified power conversion
efficiency (PCE) of 27.31\% for single cells and 21.54\% for a 766 cm2
module under International Summit on Organic Photovoltaic Stability
(ISOS) protocol ISOS-L-1 {[}9{]}. Tailoring asymmetric benzothienoindole
SAMs with electron-withdrawing substituents similarly improves packing
and energetics, delivering 25.38\% on 1 cm2 and 20.59\% on a 655.2 cm2
module {[}58{]}. At the top contact, multidentate fullerenes anchor
phenethylammonium bromide capping layers against thermal loss, retaining
98.7\% initial PCE after 2100 h of maximum power point tracking (MPPT)
under continuous illumination and 95.8\% after 1800 h at 85 °C under
ISOS-L-3 {[}59{]}, while phase-pure two-dimensional (2D) perovskite
reconstruction passivates surfaces {[}60,61{]}. Extrinsic ion migration
also accelerates failure: lithium bis(trifluoromethanesulfonyl)imide
(Li-TFSI) dopants in HTLs release mobile lithium ions that drift under
bias, whereas volatile dopants eliminate ionic residues {[}62{]}.
Surface passivation via aromatic carbonyl molecules yields 26.12\% PCE,
but only on an uncertified, sub-0.1 cm2 area (0.09 cm2) {[}63{]}.

In wide-bandgap (WBG) mixed-halide absorbers, operational instability
stems from electric-field-induced halide segregation driven by
trapped-charge accumulation {[}64{]}. Operating cells under maximum
stability point tracking (MSPT) instead of MPPT alleviates trapped
carriers, retaining approximately 90\% initial power after 700 h
{[}64{]}. Chemical stabilization targets the same thermodynamic
vulnerability: incorporating 1,4-phenylenediamine dihydriodide into
1.93-eV absorbers raises the ion migration activation barrier to
suppress phase segregation, achieving a certified PCE of 26.18\% with an
open-circuit voltage (Voc) of 3.148 V on 1.17 cm2 triple-junction cells
{[}65{]}. Rubidium incorporation similarly homogenizes vertical cesium
distribution in 1.67-eV pure-iodide perovskites to prevent demixing,
reaching an uncertified 22.51\% PCE under ISOS-L-1 {[}66{]}.

Under thermal and light fluctuations, dynamic lattice strain drives
degradation {[}67{]}: residual tensile stress relaxes via
diffusion-mediated inelastic deformation, accelerating moisture uptake
by an order of magnitude {[}68{]}. Interfacial networks and printed
silicon oxynitride barrier films counter moisture penetration, the
latter retaining 93\% ± 5\% PCE after 1000 h at 25 °C and 60\% relative
humidity (RH) under ISOS-D-1 {[}69,70,71{]}. Yet a methodological gap
obscures durability: only 4 of 169 closely read studies reported the
time to 80\% initial efficiency (T80), spanning 200 h, 720 h, 6785 h,
and 9800 h {[}5,50,72,73{]}, and only 6 specified an ISOS protocol
{[}9,11,16,17,59,66{]}. Prevailing tests apply constant 1-sun
illumination at mild temperatures, probing photocarrier collection while
failing to evaluate thermomechanical fatigue, reverse bias, or
dark-light cycles.

\begin{figure}[htbp]
\centering
\includegraphics[width=0.95\textwidth]{fig/F4_stability_evidence.pdf}
\caption{The month's operational-stability evidence in full. Left, every reported time to 80\% of initial performance. Right, the stability protocol labels used.}
\end{figure}

\subsection{7. Scale-Up and the Area
Penalty}\label{scale-up-and-the-area-penalty}

The area penalty during module scale-up stems from three coupled
physical bottlenecks: series-resistance losses across monolithic
interconnects, non-uniform crystallization in scalable coating, and
localized interfacial recombination {[}10,74,75{]}. Extrinsic
integration and resistive interconnect losses rival intrinsic material
recombination in driving the efficiency drop from discrete cells to
series-connected modules {[}74,75{]}. Hyperspectral luminescence imaging
attributes this deficit to resistive transport across interconnections
in 64 cm2 modules {[}74{]}, while electrostatic templating with
indol-3-ylethylammonium iodide spacer cations leaves a certified power
conversion efficiency (PCE) gap between 26.94\% on small-area cells and
23.21\% on a 25 cm2 monolithic module {[}75{]}. Aperture scaling
demonstrates that modules can maintain high efficiencies: a certified
27.3\% PCE (champion 27.8\%) across 625 cm2 with nonfullerene
electron-transport layers (ETLs) {[}10{]}, alongside uncertified
champion outputs of 27.12\% on 655.2 cm2 {[}4{]}, 27.31\% on 684.75 cm2
{[}8{]}, 27.31\% on 766 cm2 {[}9{]}, and 26.53\% on 818 cm2 {[}11{]}.

Transitioning from spin-coating to scalable meniscus blade coating or
vapor methods induces drying gradients and phase degradation
{[}76,77,78,79{]}. Incorporating Lewis base and zwitterion dual
additives into blade-coated tin-lead perovskite solar cells (PSCs)
yields grains over 800 nm and elevates open-circuit voltage (Voc) from
0.797 to 0.850 V, achieving 21.02\% single-junction PCE and a
two-terminal (2T) tandem champion PCE of 27.84\% (27.33\% stabilised)
{[}76{]}. For ambient fabrication, surface-confined protection
stabilizes pre-annealing crystallization, enabling air-processed
blade-coated methylammonium (MA)-free PSCs to achieve 26.14\% champion
and 25.75\% certified PCE, with 93.11\% retention after 1300 h of
continuous 1-sun maximum power point tracking (MPPT) {[}77{]}.
Vapor-phase annealing (VPA) establishes uniform thermal-humidity fields
to direct bottom-up crystallization, yielding an additive-free champion
PCE of 22.15\% {[}78{]}, whereas stabilizing precursor solutions against
aging via anion-pi interactions preserves 25.76\% efficiency after 7
days and 23.88\% after 3 months {[}79{]}. High-throughput quantification
of microstructural homogeneity is assisted by deep segmentation
frameworks such as PerovSegNet, which reaches 87.25\% mean average
precision (mAP) in scanning electron microscopy (SEM) image analysis at
25.22\% lower computational cost {[}80{]}. Replacing fullerene ETLs with
the nonfullerene conjugated polymer 2PB-T also addresses contact
degradation, retaining over 98.6\% of initial PCE after 1752 h of
continuous MPPT at 85 °C in air on small-area cells {[}10{]}.

\begin{figure}[htbp]
\centering
\includegraphics[width=0.95\textwidth]{fig/F2_area_penalty.pdf}
\caption{Efficiency against reported active or aperture area. The shaded band marks the sub-0.3 cm$^2$ laboratory-cell regime in which most champion values are measured.}
\end{figure}

\subsection{8. Research Gaps and
Outlook}\label{research-gaps-and-outlook}

The performance divergence between single junctions and tandems defines
the first gap: the best certified perovskite/silicon tandem this month
was 33.1\% power conversion efficiency (PCE) {[}1{]}, while certified
single junctions clustered near 27\%, showing that the wide-bandgap
(WBG) subcell limiting the tandem is not the one most papers optimise.
Chemical treatments claimed to homogenize mixed halides in pristine WBG
absorbers leave unresolved whether halide homogenisation suppresses
phase segregation under operating bias, as carrier gradients trigger
iodide clustering {[}49,64{]}. Operando nanoscale photoluminescence
spectroscopy during continuous maximum power point tracking (MPPT) load
under elevated operating temperatures would establish whether
compositional uniformity survives operational electric fields.

A second gap concerns buried-interface chemistry that survives thermal
cycling rather than only lifting initial efficiency. Molecules and
self-assembled monolayers (SAMs) passivate buried defects to lift
open-circuit voltage (Voc) on pristine substrates, but thermal expansion
mismatches induce cyclic shear stresses that cause delamination
{[}1,19{]}. In situ grazing-incidence wide-angle X-ray scattering with
operando transient photovoltage measurements during temperature cycling
across thermal extremes would quantify strain relaxation and resolve
whether interfacial bonding withstands thermal fatigue.

A third gap lies in ion-migration timescales measured under load rather
than in the dark. Only 4 of 169 closely read papers reported a time to
80\% of initial performance (T80) lifetime, the longest 9800 h {[}5{]},
short against a 25-year module warranty. Extracting kinetics from dark
impedance ignores how photogenerated carriers screen fields and alter
vacancy formation enthalpies under illumination {[}64,65{]}.
Intensity-modulated photovoltage spectroscopy under bias illumination
would capture true ionic drift velocities and accumulation rates under
load.

The fourth gap is whether passivation that lifts Voc on a 0.1 cm2 cell
survives slot-die coating at module area. Only 27 of 648 works addressed
scale-up at all, and the largest reported aperture was 818 cm2 {[}11{]},
orders of magnitude below a commercial module. Meniscus coating promotes
passivator aggregation, forming pinholes across large areas {[}75{]}.
Operando photoluminescence imaging synchronized with optical
reflectometry during slot-die deposition would track drying kinetics and
isolate spatial Voc losses across modules. Next month's literature would
gain immediate utility if studies paired small-cell efficiencies with
operando stability under load and reported spatial coating uniformity.

\subsection{Acknowledgements}\label{acknowledgements}

The authors would like to acknowledge the financial support from Xiamen
University Malaysia through the Xiamen University Malaysia Research Fund
(XMUMRF/2022-C10/IENG/0046), and from National Natural Science
Foundation of China (Ref. No.: 52503390).

\subsection{Declaration of Competing
Interest}\label{declaration-of-competing-interest}

The authors declare no conflict of interest.

\subsection{Data Availability
Statement}\label{data-availability-statement}

The corpus table for the month, the per-paper extraction records with
their verbatim source quotations, the statistics file underlying every
number in the text, and the gate report are provided as Supplementary
Information. Abstract text and full texts are not redistributed.

\subsection{AI Usage Declaration}\label{ai-usage-declaration}

This manuscript was prepared as part of our Research Project for making
a reliable Autonomous Researcher Agent. The agent executed the monthly
review pipeline without human intervention: it harvested the indexed
literature for the month through the OpenAlex, Semantic Scholar,
Crossref and arXiv interfaces, applied a deterministic topical filter,
ranked and selected the papers read closely, extracted every
quantitative claim into a structured record binding each number to a
verbatim quotation from that paper's own abstract, computed the
statistics and figures from those records, and drafted this text from
the resulting evidence briefs. Literature harvesting, selection,
statistics, figures and gate checking are deterministic scripts rather
than model output. Structured extraction used qwen3.8-flash through the
opencode interface; the prose was written by gemini-3.8-flash-high
through the agy interface; independent review used opus at high
reasoning effort. No figure is a model-generated image: every panel is
plotted directly from the extracted records. Every quantitative
statement in the main text resolves to the machine-generated statistics
file shipped as Supplementary Information, and every citation resolves
to a record carrying its source quotation. The author reviewed and
edited the content, verified the selection, and takes full
responsibility for the content of the published article.

\subsection{References}\label{references}

\begin{enumerate}
\def\labelenumi{\arabic{enumi}.}
\tightlist
\item
  Xu, X. et al.~Noncoplanar multidentate contacts reinforce buried
  interface for perovskite/silicon tandem solar cells. \emph{Science
  Advances} (2026). \url{https://doi.org/10.1126/sciadv.aei2945}
\item
  Wang, F. et al.~Beyond Self‐Assembled Molecules:
  Hydrogen‐Bonding‐Regulated Crystallization of Wide‐Bandgap Perovskite
  for High‐Performance Perovskite‐Silicon Tandem Photovoltaics.
  \emph{Angewandte Chemie International Edition} (2026).
  \url{https://doi.org/10.1002/anie.1676645}
\item
  Liu, Z. et al.~Homogenizing Sn--Pb Distribution Through A‐Site MA
  Cations for Efficient All‐Perovskite Tandem Solar Cells.
  \emph{Advanced Materials} (2026).
  \url{https://doi.org/10.1002/adma.74078}
\item
  Yang, R. et al.~Soft supermolecule stabilized buried interface for
  high-performance inverted perovskite solar cells and modules.
  \emph{Nature Communications} (2026).
  \url{https://doi.org/10.1038/s41467-026-74018-8}
\item
  Zheng, W. et al.~Intermediate‐Site Anchoring Ligands Enable Robust
  Nonlayered Interfacial Passivation for Efficient and Stable
  Air‐Processed Perovskite Solar Cells. \emph{Advanced Materials}
  (2026). \url{https://doi.org/10.1002/adma.74243}
\item
  Lp, H. et al.~Enhancing Sn-Pb perovskite homogeneity via thioether
  coordination for efficient and stable all-perovskite tandem solar
  cells. \emph{Science Advances} (2026).
  \url{https://doi.org/10.1126/sciadv.aeb8790}
\item
  Xiong, Y. et al.~Radiation-resilient monolithic wide-bandgap
  perovskite/p-type heterojunction silicon tandem solar cells for space
  photovoltaics. \emph{Science Advances} (2026).
  \url{https://doi.org/10.1126/sciadv.aef6600}
\item
  Li, B.E. et al.~Unlocking 27.3\% Perovskite Solar Cells and
  Ultra‐Stable Flexible Modules via Multidentate Molecular‐Mediated
  Buried‐Interface Reconstruction. \emph{Advanced Materials} (2026).
  \url{https://doi.org/10.1002/adma.73807}
\item
  Zhang, Z.Y. et al.~Unlocking 27.3\% perovskite photovoltaics by
  interface-locked dual-molecule contact. \emph{Science Advances}
  (2026). \url{https://doi.org/10.1126/sciadv.aeg1456}
\item
  Gao, D. et al.~n-Type polymer layers enable efficient, scalable, and
  thermally stable perovskite solar modules. \emph{Science} (2026).
  \url{https://doi.org/10.1126/science.aed8175}
\item
  Qian, C. et al.~Bidirectional Chemo-Mechanical Interface Stabilization
  in Perovskite Solar Cells. \emph{Journal of the American Chemical
  Society} (2026). \url{https://doi.org/10.1021/jacs.6c06932}
\item
  Li, R. et al.~Recent Advancing Progress of Inverted Perovskite Solar
  Cells Through Functional Buried Interface Engineering. \emph{Advanced
  Functional Materials} (2026). \url{https://doi.org/10.1002/adfm.77486}
\item
  Li, M. et al.~Improved Mechanical Adhesion Enabled by Bilayer Hole
  Transport Layers in Perovskite Solar Cells. \emph{ACS Materials
  Letters} (2026).
  \url{https://doi.org/10.1021/acsmaterialslett.6c00479}
\item
  Wu, H. et al.~A Multidentate Coordination Complex‐Based Self‐Assembled
  Monolayer Stabilizes Perovskite Interface. \emph{Angewandte Chemie
  International Edition} (2026).
  \url{https://doi.org/10.1002/anie.8080769}
\item
  Duan, Y. et al.~Bisphosphonate Self‐Assembled Molecules Enhance
  Anchoring Strength and Perovskite Passivation Toward High‐Efficiency
  Perovskite/Silicon Tandem Solar Cells. \emph{Small} (2026).
  \url{https://doi.org/10.1002/smll.74999}
\item
  Qin, Z. et al.~Tetraphosphorylated phthalocyanine--based
  self-assembled monolayer stabilizes perovskite photovoltaics.
  \emph{Science Advances} (2026).
  \url{https://doi.org/10.1126/sciadv.aee4175}
\item
  Wang, J. et al.~Suppression of SAM desorption by engineering
  polyoxometalate interface for inverted perovskite solar cells.
  \emph{Polyoxometalates} (2026).
  \url{https://doi.org/10.26599/pom.2026.9140147}
\item
  Chen, X. et al.~Molecular Lubricant Mitigates Self‐Assembled Monolayer
  Aggregation in Perovskite Photovoltaics. \emph{Angewandte Chemie
  International Edition} (2026).
  \url{https://doi.org/10.1002/anie.7090517}
\item
  Zhou, B. et al.~Proton Transfer Complex-Confined Self-Assembled
  Monolayer for Thermally Stable Inverted Perovskite Solar Cells.
  \emph{ACS Energy Letters} (2026).
  \url{https://doi.org/10.1021/acsenergylett.6c01489}
\item
  Wan, X. et al.~Strengthened Charge‐Selective Perovskite/SAM
  Heterocontact for Efficient and Durable Perovskite Solar Cells via
  Interfacial Reconfiguration. \emph{Angewandte Chemie International
  Edition} (2026). \url{https://doi.org/10.1002/anie.1018855}
\item
  Ling, X. et al.~In situ Induced 2D Perovskite Enhances Buried
  Interface Passivation for Ambient‐Processed Efficient Perovskite Solar
  Cells. \emph{Advanced Functional Materials} (2026).
  \url{https://doi.org/10.1002/adfm.76848}
\item
  Xu, Y. et al.~Eliminating Buried Interface Voids for High‐Efficiency
  and Stable Perovskite Solar Cells. \emph{Small} (2026).
  \url{https://doi.org/10.1002/smll.74604}
\item
  Wang, X.F. et al.~Interfacial Ionic Redistribution Induced by
  Sulfonate‐Anchored Dipotassium Salts Enables Efficient Buried
  Interface Engineering in Perovskite Solar Cells. \emph{Advanced
  Functional Materials} (2026). \url{https://doi.org/10.1002/adfm.77175}
\item
  Su, H. et al.~Synergistic Interface Modulation of Hole Transport for
  Efficient, Stable Inverted Perovskite Solar Cells. \emph{Advanced
  Functional Materials} (2026). \url{https://doi.org/10.1002/adfm.76878}
\item
  Chen, X. et al.~Perylenocarbazole‐Based Polycyclic Aromatic
  Self‐Assembled Monolayers with Tailored Electrostatic Potentials for
  High‐Performance Organic and Perovskite Solar Cells. \emph{Angewandte
  Chemie International Edition} (2026).
  \url{https://doi.org/10.1002/anie.4031162}
\item
  Ye, Y. et al.~Interfacial Regulation by a Multifunctional
  Self‐Assembled Material for High‐Performance Inverted Perovskite Solar
  Cells. \emph{Advanced Energy Materials} (2026).
  \url{https://doi.org/10.1002/aenm.71282}
\item
  Ma, P. et al.~Ternary Cooperative Interface With Electronegative
  Phosphine Bridges Enables Efficient and Stable Inverted Perovskite
  Solar Cells. \emph{Advanced Materials} (2026).
  \url{https://doi.org/10.1002/adma.74416}
\item
  Dutta, R. et al.~Temperature‐Dependent Performance Analysis of
  Formamidinium‐Based Perovskite Solar Cells: A Theoretical Study.
  \emph{Energy Technology} (2026).
  \url{https://doi.org/10.1002/ente.70590}
\item
  Jubu, P.R. et al.~Interface-engineered HTL-free mixed-cation
  perovskite solar cells with ZnO/BaSnO3 bilayer electron transport
  architecture: A SCAPS-1D investigation toward near-ideal photovoltaic
  performance. \emph{Hybrid Advances} (2026).
  \url{https://doi.org/10.1016/j.hybadv.2026.100721}
\item
  Wong, C. et al.~Rational Tailoring of Hole‐Selective Self‐Assembly
  Monolayers Based on Sulfur‐Containing Heterocycles for
  High‐Performance Perovskite Solar Cells. \emph{Angewandte Chemie
  International Edition} (2026).
  \url{https://doi.org/10.1002/anie.1864800}
\item
  Shang, X. et al.~A Multisite Molecular Salt Enables Surface
  Passivation and Energetics Management for Efficient Inverted
  Perovskite Solar Cells. \emph{Advanced Energy Materials} (2026).
  \url{https://doi.org/10.1002/aenm.71305}
\item
  Tan, Y. et al.~Synergistic Interfacial Dangling‐Bond Passivation for
  Enhanced Moisture Resistance in Perovskite Photovoltaics.
  \emph{Advanced Materials} (2026).
  \url{https://doi.org/10.1002/adma.73950}
\item
  Saeed, A. et al.~Amphoteric Anion‐Mediated Crystallization and
  Universal Defect Passivation for High‐Efficiency and Stable Perovskite
  Solar Cells. \emph{Advanced Functional Materials} (2026).
  \url{https://doi.org/10.1002/adfm.77378}
\item
  Chen, R. et al.~Reconstructing the perovskite surface with a
  homogeneous electric potential by reducing localized charge for
  efficient and durable perovskite solar cells. \emph{Energy Materials}
  (2026). \url{https://doi.org/10.20517/energymater.2026.102}
\item
  Jang, W.J. et al.~Conjugated anion modulation for extended defect
  passivation at buried interfaces in perovskite solar cells.
  \emph{Advanced Composites and Hybrid Materials} (2026).
  \url{https://doi.org/10.1007/s42114-026-01958-6}
\item
  Fan, R. et al.~Plasma surface engineering for efficient and stable
  perovskite solar cells and modules. \emph{Science} (2026).
  \url{https://doi.org/10.1126/science.aeg1730}
\item
  Deng, Z. et al.~Precise Iodide‐Trap Strategy for Photothermally Stable
  Spiro‐OMeTAD‐Based Perovskite Solar Cells. \emph{Angewandte Chemie
  International Edition} (2026).
  \url{https://doi.org/10.1002/anie.2924886}
\item
  Miao, Z. et al.~A unified ligand-dimensional design to halt cation
  migration in perovskite photovoltaics. \emph{Science Advances} (2026).
  \url{https://doi.org/10.1126/sciadv.aed6327}
\item
  Du, X. et al.~Enhancing Interfacial Stability of FAPbI 3 Perovskite
  Solar Cells by Introducing 2D/3D Heterojunction Based on Functional
  Imidazolium Cations. \emph{Advanced Functional Materials} (2026).
  \url{https://doi.org/10.1002/adfm.77015}
\item
  Cheng, J. et al.~Modulating Growth Kinetics and Defect Passivation in
  Ambient‐Air Evaporation‐Spray Deposition Enables Efficient and Stable
  Rigid/Flexible Perovskite Solar Cells. \emph{EcoEnergy} (2026).
  \url{https://doi.org/10.1002/ece2.70102}
\item
  LIU, M. et al.~In‐Situ Constructed Cations for 2D/3D Perovskite
  Heterostructure for Stable and Efficient Photovoltaics. \emph{Advanced
  Science} (2026). \url{https://doi.org/10.1002/advs.76508}
\item
  Zhu, W. et al.~Thermodynamically Matched Molten‐Salt‐Assisted Phase
  Engineering of 2D/3D Perovskite Heterostructures for Efficient and
  Stable Solar Cells. \emph{Advanced Functional Materials} (2026).
  \url{https://doi.org/10.1002/adfm.76974}
\item
  Lv, D. et al.~Dual‐Site‐Binding Ligand‐Driven (100) Orientation for
  Efficient and Stable Perovskite Solar Cells. \emph{Advanced Materials}
  (2026). \url{https://doi.org/10.1002/adma.74033}
\item
  Wang, D. et al.~Multi‐Bonding Molecular Interface Bridges Enable
  Synergistic Defect Passivation, Strain Relaxation, and Accelerated
  Electron Extraction in Perovskite Solar Cells. \emph{Advanced
  Functional Materials} (2026). \url{https://doi.org/10.1002/adfm.77373}
\item
  Yang, J. et al.~Dynamic Strain Regulation Via Photoresponsive
  Fullerenes Enables High‐Performance and UV‐Robust Perovskite Solar
  Cells. \emph{Advanced Materials} (2026).
  \url{https://doi.org/10.1002/adma.74220}
\item
  Yang, S. et al.~Low-temperature liquid-assisted vibrational
  stress-relief strategy enables stable and efficient perovskite solar
  cells. \emph{Nature Communications} (2026).
  \url{https://doi.org/10.1038/s41467-026-75629-x}
\item
  Asare, G.K. et al.~Redox‐Buffering PEAI‐CeO x Passivation Suppresses
  Thermo‐Mechanical Degradation in Perovskite Solar Cells Under Rapid
  Thermal Cycling. \emph{Advanced Energy Materials} (2026).
  \url{https://doi.org/10.1002/aenm.71312}
\item
  Mardegan, L. et al.~Materials and Device Engineering for Efficient,
  Stable, and Scalable Monolithic Perovskite/Silicon Tandem
  Photovoltaics. \emph{Chemical Reviews} (2026).
  \url{https://doi.org/10.1021/acs.chemrev.5c01014}
\item
  Durant, B.K. et al.~In-Operando Thermal Cycling of All-Perovskite
  Tandem Solar Cells. \emph{ACS Energy Letters} (2026).
  \url{https://doi.org/10.1021/acsenergylett.6c00992}
\item
  Liu, P. et al.~Sr/Pb Alloying Stabilizes Wide‐Bandgap Perovskites for
  Durable Fully‐Textured Perovskite/Silicon Tandem Solar Cells.
  \emph{Advanced Energy Materials} (2026).
  \url{https://doi.org/10.1002/aenm.71317}
\item
  Chen, Y. et al.~Molecular‐Extrusion‐Driven Halogen Homogenization for
  Efficient Perovskite‐Silicon Tandem Solar Cells. \emph{Angewandte
  Chemie International Edition} (2026).
  \url{https://doi.org/10.1002/anie.5530609}
\item
  Zhou, Y. et al.~Conductive Polymer Bridges Mediate High‐Performance
  Flexible Perovskite--Organic Tandem Solar Cells. \emph{Advanced
  Materials} (2026). \url{https://doi.org/10.1002/adma.74183}
\item
  Liu, L. et al.~Artificial Intelligence-Guided Cosolvent Design for
  High-Performance Perovskite/Silicon Tandem Solar Cells.
  \emph{Nano-Micro Letters} (2026).
  \url{https://doi.org/10.1007/s40820-026-02291-9}
\item
  Dangudubiyyam, U.K. et al.~Eco‐Friendly Solvent System for Inkjet
  Deposition of Wide Bandgap Perovskite Solar Cells Enabling Tandem
  Integration. \emph{Advanced Science} (2026).
  \url{https://doi.org/10.1002/advs.76463}
\item
  Yun, Y. et al.~Triple-junction all-perovskite solar cells with
  self-assembling hole contacts in all subcells. \emph{Joule} (2026).
  \url{https://doi.org/10.1016/j.joule.2026.102575}
\item
  Rehan, M. et al.~Efficient Flexible Monolithic Perovskite/CIGS Tandem
  Solar Cell by Using Two‐Step Hybrid Deposition Approach.
  \emph{Advanced Energy Materials} (2026).
  \url{https://doi.org/10.1002/aenm.71267}
\item
  Chalkias, D.A. et al.~Depth‐Resolved Photo‐Degradation of Perovskite
  Layer in Solar Cells Stressed Under ISOS Light‐Aging Protocol
  Conditions: The Role of Substrate, Electron‐Transport‐Layer and
  Perovskite Deposition Method. \emph{Solar RRL} (2026).
  \url{https://doi.org/10.1002/solr.70421}
\item
  Zhong, S. et al.~Molecular Tailoring of Self‐Assembled Monolayers
  Derived From Asymmetric Benzothienoindole for Efficient and Stable
  Inverted Perovskite Solar Cells and Modules. \emph{Angewandte Chemie
  International Edition} (2026).
  \url{https://doi.org/10.1002/anie.1678852}
\item
  Huang, J. et al.~Multidentate Fullerene Boosts the Thermal Stability
  of Perovskite Solar Cells. \emph{ACS Energy Letters} (2026).
  \url{https://doi.org/10.1021/acsenergylett.6c01517}
\item
  Dai, K. et al.~Thermally Independent Interfacial Reconstruction for
  Phase‐Pure n = 1 2D Perovskites With Mixed‐Orientation Surface
  Architecture. \emph{Advanced Materials} (2026).
  \url{https://doi.org/10.1002/adma.74247}
\item
  You, S. et al.~Boosting Carrier Transport in Carbon‐Based
  Hole‐Transport‐Layer‐Free Perovskite Solar Cells via 2D/3D Interfacial
  Structural Configuration. \emph{Solar RRL} (2026).
  \url{https://doi.org/10.1002/solr.70447}
\item
  Zhao, J. et al.~Eliminating Residue‐Induced Degradation: A Volatile
  Dopant Strategy for High Performance Perovskite Photovoltaics.
  \emph{Small} (2026). \url{https://doi.org/10.1002/smll.74289}
\item
  Yang, G. et al.~Central Core Engineering of Aromatic Carbonyl
  Molecules Enables Highly Efficient and Stable Perovskite Solar Cells.
  \emph{Angewandte Chemie International Edition} (2026).
  \url{https://doi.org/10.1002/anie.5216345}
\item
  Kim, S. et al.~Maximum Stability Point Tracking Stabilizes
  Wide‐Bandgap Mixed‐Halide Perovskite Solar Cells. \emph{Advanced
  Energy Materials} (2026). \url{https://doi.org/10.1002/aenm.71358}
\item
  Kong, W. et al.~Stabilizing 1.93-eV ultrawide-bandgap perovskites for
  efficient triple-junction solar cells. \emph{Nature Communications}
  (2026). \url{https://doi.org/10.1038/s41467-026-75305-0}
\item
  Zhou, X. et al.~Homogenizing Cesium Distribution via Rubidium
  Incorporation Enables Pure-Iodide 1.67 eV Bandgap Perovskite Solar
  Cells with Efficiency Exceeding 22\%. \emph{ACS Nano} (2026).
  \url{https://doi.org/10.1021/acsnano.6c04084}
\item
  Fasti, R. et al.~Dynamic Stress and Strain Evolution and Engineering
  Strategies to Enhance Operational Durability in Perovskite Solar
  Cells. \emph{ACS Energy Letters} (2026).
  \url{https://doi.org/10.1021/acsenergylett.6c01837}
\item
  Guo, B. et al.~A Composition-Dependent Stress--Diffusion Framework
  Identifies Humidity-Tolerant Processing Windows in Hybrid Perovskite
  Solar Cells. \emph{ChemRxiv} (2026).
  \url{https://doi.org/10.26434/chemrxiv-2025-128tn/v2}
\item
  Li, B. et al.~Inkjet‐Printed SiO x N y Barrier Film for Encapsulation
  of Perovskite Solar Cells. \emph{Small} (2026).
  \url{https://doi.org/10.1002/smll.74816}
\item
  Albab, M.F. et al.~Air-Processed and Water-Stable Perovskite Solar
  Cells Enabled by a Fishing-Net-Inspired Interfacial Network.
  \emph{Nano-Micro Letters} (2026).
  \url{https://doi.org/10.1007/s40820-026-02295-5}
\item
  Gao, B. et al.~Spatiotemporally homogeneous crystallization for
  ambient scalable perovskite photovoltaics. \emph{Science} (2026).
  \url{https://doi.org/10.1126/science.aef1969}
\item
  Ahmad, N. et al.~Synergistic Interface and Band Engineering for Sn‐Pb
  Perovskites. \emph{Small} (2026).
  \url{https://doi.org/10.1002/smll.74531}
\item
  Ni, Q. et al.~Protonated Cyclohexylamine as an Effective Passivation
  Layer for High-Performance Perovskite Solar Cells. \emph{ACS Applied
  Energy Materials} (2026). \url{https://doi.org/10.1021/acsaem.6c01637}
\item
  Harada, N. et al.~Quantitative Hyperspectral Luminescence Imaging
  Reveals Upscaling Losses in Perovskite Solar Modules. \emph{Solar RRL}
  (2026). \url{https://doi.org/10.1002/solr.70375}
\item
  Ki, T. et al.~Asymmetric Lamellar Templating of the Perovskite/C 60
  Interface for Scalable Inverted Perovskite Photovoltaics.
  \emph{Advanced Materials} (2026).
  \url{https://doi.org/10.1002/adma.74316}
\item
  Wu, Y. et al.~Lewis Base--Zwitterion Dual Additives Enable Efficient
  Blade‐Coated Tin‐lead Perovskite Solar Cells and Tandems.
  \emph{Advanced Energy Materials} (2026).
  \url{https://doi.org/10.1002/aenm.71336}
\item
  Duan, L. et al.~Surface-confined protection stabilizes pre-annealing
  crystallization for ambient blade-coated perovskites. \emph{Nature
  Communications} (2026).
  \url{https://doi.org/10.1038/s41467-026-75053-1}
\item
  Xu, C. et al.~Vapor-phase Annealing toward Homogeneous Bottom-Up
  Crystallization in Perovskite Solar Cells. \emph{ACS Applied Materials
  \& Interfaces} (2026). \url{https://doi.org/10.1021/acsami.6c07418}
\item
  Sun, T. et al.~Stabilizing Precursor Solutions via Anion‐π Interaction
  for Robust Production of Perovskite Solar Cells. \emph{Advanced
  Functional Materials} (2026). \url{https://doi.org/10.1002/adfm.77211}
\item
  Wang, L. et al.~Automated and Scalable SEM Image Analysis of
  Perovskite Solar Cell Materials via a Deep Segmentation Framework.
  \emph{npj Computational Materials} (2026).
  \url{https://doi.org/10.1038/s41524-026-02211-w}
\end{enumerate}

\end{document}
